% SEGMENT OLP-0013-S01
% Part: sets-functions-relations
% Chapter: relations
% Section: reflections
%
\documentclass[../../../include/open-logic-section]{subfiles}

\begin{document}

\olfileid{sfr}{rel}{ref}
\olsection{மெய்யியல் சிந்தனைகள்}

\olref[set]{sec} இல், தொடர்புகளைச் சில வகைக் கணங்களாக வரையறுத்தோம்.
சற்றே நின்று ஒரு சிறிய மெய்யியல் கேள்வியைக் கேட்போம்:
அத்தகைய வரையறை என்ன \emph{செய்கிறது}?
மெய்ப்பொருளியல் சார்ந்த ஒன்றாதல் உண்மைகள் சிலவற்றை நாம்
\emph{கண்டுபிடித்துவிட்டோம்} எனக் கூற விரும்புவது மிகவும் ஐயத்திற்குரியது.
எடுத்துக்காட்டாக, $\Nat$ இன் மீதான வரிசைத் தொடர்பு,
\olref[set]{sec} இல் வரையறுத்த
$R=\Setabs{\tuple{n,m}}{n, m \in \Nat\text{ மற்றும் } n < m}$
என்னும் கணமே \emph{எனத் தெரியவந்தது} என்று கூறுவதை எடுத்துக்கொள்வோம்.
இதனை ஐயுறுவதற்கு மூன்று காரணங்கள் உள்ளன.

% SEGMENT OLP-0013-S02
முதலாவது: \olref[set][pai]{wienerkuratowski} இல்
$\tuple{a, b} = \{\{a\}, \{a, b\}\}$ என வரையறுத்தோம்.
இதற்குப் பதிலாக $\lVert a, b\rVert = \{\{b\}, \{a, b\}\} = \tuple{b,a}$
என்னும் வரையறையைக் கருதுக. $a \neq b$ எனில்,
$\tuple{a, b} \neq \lVert a,b\rVert$.
ஆனால் வரிசைப்படுத்தப்பட்ட ஜோடிக்கான நமது வரையறையாக
$\tuple{a,b}$ என்பதற்குப் பதிலாக $\lVert a,b\rVert$ என்பதையும்
அதே அளவு பொருத்தத்துடன் எடுத்திருக்கலாம்.
இரண்டு வரையறைகளும் ஒரே அளவுக்கு நன்றாகச் செயல்பட்டிருக்கும்.
எனவே இயல் எண்களின் மீதான வரிசைத் தொடர்பாக ``இருப்பதற்கு''
இப்போது சமமான தகுதியுள்ள இரண்டு வாய்ப்புகள் உள்ளன. அவை:
\begin{align*}
		R &= \Setabs{\tuple{n,m}}{n, m \in \Nat \text{ மற்றும் }n < m}\\
		S &= \Setabs{\lVert n,m\rVert}{n, m \in \Nat \text{ மற்றும் }n < m}.
\end{align*}
உறுப்புசார் சமத்துவத்தின்படி $R \neq S$ என்பதால்,
இவை \emph{இரண்டுமே} $\Nat$ இன் மீதான வரிசைத் தொடர்புடன் ஒன்றாக இருக்க
முடியாது என்பது தெளிவு. ஆனால் $S$ என்பதற்குப் பதிலாக $R$
(அல்லது மறுதிசையில்) உண்மையில் வரிசைத் தொடர்பாக \emph{இருக்கிறது}
எனக் கூறுவது தன்னிச்சையானதாகவும், எனவே சற்றே சங்கடமானதாகவும் இருக்கும்.
(கணக்கோட்பாட்டுச் சுருக்கவாதத்திற்கு எதிராக பெனசெராஃப்
\citeyear{Benacerraf1965} இல் புகழ்பெறச் செய்த வாதத்தின் மிக எளிய
எடுத்துக்காட்டு இது. இதனைப் பலமுறை மீண்டும் காண்போம்.)

% SEGMENT OLP-0013-S03
இரண்டாவது: \emph{ஒவ்வொரு} தொடர்பையும் ஒரு கணத்துடன் ஒன்றாகக் கருத வேண்டும்
என்று நினைத்தால், கண-உறுப்பாதல் தொடர்பான $\in$ என்பதே ஒரு குறிப்பிட்ட
கணமாக இருக்க வேண்டும். உண்மையில், அது
$\Setabs{\tuple{x,y}}{x \in y}$ என்ற கணமாக இருக்க வேண்டும்.
ஆனால் இந்தக் கணம் இருக்கிறதா? ரஸ்ஸலின் முரண்பாட்டைக் கருத்தில் கொண்டால்,
இத்தகைய கணம் இருக்கிறது என்பது சாதாரணமான கூற்றல்ல.
உண்மையில், \oliflabeldef{cumul:::part}{\olref[cumul][][]{part} இல்
நாம் உருவாக்கும் கணக்கோட்பாடு, இந்தக் கணம் இருப்பதை
\emph{மறுக்கும்}.\footnote{முன்னோக்கிச் சென்று அதன் காரணத்தைச் சொல்வோம்.
முரண்பாட்டின் மூலம் நிறுவுவதற்காக,
$I = \Setabs{\tuple{x,y}}{x \in y}$ இருப்பதாகக் கொள்வோம்.
அப்போது $\bigcup \bigcup I$ அனைத்துக் கணமாகும்;
இது \olref[sfr][z][sep]{thm:NoUniversalSet} க்கு முரணாகும்.}}{கணக்கோட்பாட்டை
ஓர் அடிக்கோள் கோட்பாடாகத் துல்லியமாக உருவாக்க முடியும்;
அக்கோட்பாடு இந்தக் கணம் இருப்பதை மறுக்கும்.}
எனவே சில தொடர்புகளைக் கணங்களாகக் கருத முடிந்தாலும்,
கண-உறுப்பாதல் தொடர்பு ஒரு தனிச்சிறப்பு நிலையாக இருக்க வேண்டியிருக்கும்.

% SEGMENT OLP-0013-S04
மூன்றாவது: தொடர்புகளைக் கணங்களுடன் ``ஒன்றாகக் கருதும்போது'',
$\tuple{x,y} \in R$ என்பதற்குப் பதிலாக $Rxy$ என எழுதிக்கொள்ளலாம் என்றோம்.
``$\in$'' என்னும் உறுப்பாதல் தொடர்பை ஒரு \emph{பயனிலையாகக்} கருதினால்,
இது சரியே. ஆனால் ``$\in$'' என்பது ஒரு குறிப்பிட்ட வகைக் கணத்தைக் குறிக்கிறது
என்று நினைத்தால், ``$\tuple{x,y} \in R$'' என்னும் வெளிப்பாடு,
கணங்களைக் குறிக்கும் மூன்று தனிப்பொருட் சொற்களை மட்டுமே கொண்டதாகிவிடும்:
``$\tuple{x,y}$'', ``$\in$'', ``$R$'' ஆகியவையே அவை.
இத்தகைய பெயர்ப்பட்டியல் ஒரு கூற்றை வெளிப்படுத்துவதில்,
``அந்தக் கோப்பை பேனாதாங்கி அந்த மேசை'' என்னும் பொருளற்ற சரத்தைவிட
மேம்பட்டதல்ல. மீண்டும், சில தொடர்புகளைக் கணங்களாகக் கருத முடிந்தாலும்,
கண-உறுப்பாதல் தொடர்பு தனிச்சிறப்பு நிலையாக இருக்க வேண்டும்.
(ஃப்ரேகேயின் கருத்து \emph{குதிரை} முரண்பாட்டின் ஓர் எளிய வடிவத்தையும்,
விட்கன்ஸ்டைன் ஒருமுறை ரஸ்ஸலுக்கு எதிராக எழுப்பிய புகழ்பெற்ற
மறுப்பையும் இது ஒன்றிணைக்கிறது.)

% SEGMENT OLP-0013-S05
ஆக, இதிலிருந்து நாம் பெறுவது என்ன?
எண்களின் மீதான தொடர்புகள் கணங்களாகும் எனக் கூறுவதில்
\emph{தவறு} எதுவும் இல்லை. அந்தக் கருத்து எந்த நோக்கில் சொல்லப்படுகிறது
என்பதை மட்டும் புரிந்துகொள்ள வேண்டும்.
மெய்ப்பொருளியல் சார்ந்த ஒன்றாதல் உண்மையை நாம் கூறவில்லை.
சில சூழல்களில், சில தொடர்புகளைச் சில கணங்களாகக்
\emph{கருதலாம்}, அவ்வாறே கருதுவோம் என்பதையே குறிப்பிடுகிறோம்.

\end{document}
